\chapter{Méthodologie de travail}

Ce chapitre présente la méthodologie de travail suivie pour la réalisation du projet :
l'analyse du besoin (exigences fonctionnelles et non fonctionnelles), le cahier des
charges qui a cadré le stage, ainsi que les outils et technologies retenus.

\section{Vue analytique sur le besoin}

Avant le déploiement de NetSentry, le département IT du Mövenpick Hôtel de Tanger ne
disposait d'aucun inventaire fiable des équipements connectés au réseau local : postes
de travail, imprimantes, caméras, bornes Wi-Fi et appareils invités se multiplient sans
qu'aucun outil ne permette de savoir, à un instant donné, combien d'appareils sont
actifs, ni d'identifier une connexion inhabituelle ou non autorisée.

L'analyse du besoin a permis de dégager trois exigences fonctionnelles principales :
découvrir automatiquement les équipements présents sur le réseau et les identifier
(adresse IP, adresse MAC, fabricant, ports ouverts) ; conserver un historique des
détections afin de suivre l'apparition et la disparition des appareils dans le temps ;
et signaler tout appareil détecté pour la première fois, afin que l'administrateur IT
puisse le vérifier et, le cas échéant, le classer comme connu.

Sur le plan non fonctionnel, la solution devait rester simple d'utilisation pour du
personnel non spécialisé en réseau, s'exécuter de manière autonome (scans planifiés)
sans intervention manuelle systématique, et ne transmettre aucune donnée du réseau
interne vers des services externes. Ces contraintes ont directement guidé les choix
d'architecture et de conception détaillés dans les sections suivantes.

\section{Cahier des charges}

Le cadrage du projet a été formalisé dans un cahier des charges remis en début de
stage, définissant le contexte, les objectifs, le périmètre fonctionnel et non
fonctionnel, ainsi que les contraintes et livrables attendus.

\subsection{Livrables attendus}

\begin{itemize}[leftmargin=1.2cm]
    \item Code source de l'application (backend, frontend, scripts de scan).
    \item Base de données structurée contenant l'historique des équipements détectés.
    \item Tableau de bord fonctionnel de visualisation.
    \item Image(s) Docker et instructions de déploiement.
    \item Documentation technique et rapport de stage détaillant la démarche, les choix
    techniques et les résultats obtenus.
\end{itemize}

\subsection{Contraintes}

Le projet devait être réalisé dans un délai de deux mois, avec un accès limité aux
systèmes métiers de l'hôtel (gérés par un prestataire externe). Le scan réseau ne
pouvait être exécuté que sur le réseau interne de l'hôtel, avec l'autorisation du
département IT, et la solution devait privilégier une base de données relationnelle
(SQL) en cohérence avec les outils disponibles côté hôtel.

\section{Outils et Technologies}

\subsection{Outils employés}

Les outils prévus pour ce projet sont : Node.js et Express pour le backend, une base de
données SQL pour le stockage des équipements et de leur historique, Nmap pour la
découverte réseau, ainsi que Docker pour la conteneurisation de l'application.

\subsection{Environnement de travail}

Le développement s'effectue en JavaScript (Node.js), avec un environnement de travail
basé sur Visual Studio Code et un système de gestion de versions Git.

\subsection{Synthèse des technologies retenues}

Le tableau~\ref{tab:technos} récapitule l'ensemble des choix technologiques retenus pour
chacune des couches de l'architecture NetSentry.

\begin{table}[H]
    \centering
    \caption{Synthèse des technologies retenues par couche}
    \label{tab:technos}
    \begin{tabular}{@{}>{\raggedright\arraybackslash}p{3.6cm} >{\raggedright\arraybackslash}p{4.4cm} >{\raggedright\arraybackslash}p{6.6cm}@{}}
        \toprule
        \textbf{Couche} & \textbf{Technologie(s)} & \textbf{Rôle principal} \\
        \midrule
        Présentation & React (Vite) & Tableau de bord web \\
        Logique métier (back-end) & Node.js / Express & API REST, orchestration, authentification (JWT) \\
        Scan réseau & Nmap & Découverte des hôtes (ARP) et énumération des ports \\
        Base de données & PostgreSQL & Équipements, historique des scans, alertes \\
        Conteneurisation & Docker / Docker Compose & Empaquetage et déploiement de la solution \\
        Gestion de version & Git / GitHub & Suivi des modifications du code source \\
        \bottomrule
    \end{tabular}
\end{table}
